\documentclass[minitoc, hidelinks, headlogo, framedcover]{tfeEPCC}
\usepackage[spanish, es-tabla]{babel} % Comentar si se va a realizar en inglés
%\usepackage{parskip} % Quitar comentario si se quieren distinguir los párrafos con saltos y quitar las sangrías

\usepackage{lipsum} % Este paquete es para poner texto genérico en los capítulos, quitar cuando ya se hayan redactado
\usepackage{graphicx}
\usepackage{xcolor}
\usepackage{pifont}
\usepackage{etoolbox}
\usepackage{caption}
\usepackage{listings}

\definecolor{LightBlue}{RGB}{200,230,255} % o usa blue!10 para algo más sencillo
\definecolor{mybeige}{RGB}{245, 245, 200}

\lstset{
    inputencoding=utf8,
    language=Python,
    basicstyle=\ttfamily\small,
    backgroundcolor=\color{mybeige},
    keywordstyle=\color{blue},
    commentstyle=\color{lightgray},
    stringstyle=\color{red},
    frame=single,
    rulecolor=\color{black},
    showstringspaces=false,
    numbers=left,
    numberstyle=\tiny,
    numbersep=5pt,
    tabsize=4,
    aboveskip=1em,                     % Espacio antes del bloque
    belowskip=2em,                     % Espacio después del bloque
    breaklines=true,
    breakatwhitespace=false,
    escapeinside={(*@}{@*)},
    literate={ñ}{{\~n}}1
             {á}{{\'a}}1
             {é}{{\'e}}1
             {í}{{\'i}}1
             {ó}{{\'o}}1
             {ú}{{\'u}}1
             {Á}{{\'A}}1
             {É}{{\'E}}1
             {Í}{{\'I}}1
             {Ó}{{\'O}}1
             {Ú}{{\'U}}1
             {Ñ}{{\~N}}1
}

\newcommand{\cmark}{\textcolor{green!60!black}{\ding{51}}}
\newcommand{\xmark}{\textcolor{red}{\ding{55}}}
\newcommand{\starone}{\textcolor{orange!80!black}{\ding{72}}\textcolor{lightgray}{\ding{72}\ding{72}}}
\newcommand{\startwo}{\textcolor{orange!80!black}{\ding{72}\ding{72}}\textcolor{lightgray}{\ding{72}}}
\newcommand{\starthree}{\textcolor{orange!80!black}{\ding{72}\ding{72}\ding{72}}}

\newenvironment{parskipenv}{%
  \parskip=10pt  % espacio entre párrafos
}{}

% Rellenar con los datos del estudiante y el trabajo
\estudiante{David Gordillo Burrero}{Grado}{Ingeniería Informática en Ingeniería del Software}
\title{Desarrollo de una herramienta para la automatización de tareas de reconocimiento en auditorías de seguridad}
\keywords{Seguridad\and Hacking Ético\and Auditoría\and Automatización\and Herramienta}

\begin{document}
\sloppy 
\pagebreak 

%%%%%%%%%%%%% PORTADA %%%%%%%%%%%%%
% Poner convocatoria y año de presentación
\portadaTFE{Convocatoria, 2025}

\pagestyle{empty}


%%%%%%%%%% CONTRAPORTADA %%%%%%%%%%
% Campos, separar con comas si se ponen varios: {
% studentGender
% tutor (obligatorio)
% tutorGender
% cotutor
% cotutorGender }
% Especificar un co-tutor lo añade a la contraportada
% Poner la letra f en cualquier campo de género añade una 'a' al título correspondiente (por ejemplo, studentGender = f hace que se muestre Autora)
\contraportadaTFE{tutor = Andrés Caro Lindo}

%%%% CAPÍTULOS INTRODUCTORIOS %%%%
\pagenumbering{gobble}
%\chapter*{Agradecimientos} % Apartado opcional, descomentar si se quiere añadir
% Dedicatorias del trabajo, pueden ser en español o inglés dependiendo de lo que se considere más apropiado

\chapter*{Resumen}
En el presente trabajo se ha desarrollado una herramienta orientada a la automatización de tareas fundamentales en auditorías de seguridad, con un enfoque específico en la fase de reconocimiento. Esta etapa, crítica en cualquier proceso de pentesting, abarca actividades como la recolección de información pasiva y activa, el escaneo de puertos y servicios, la enumeración de subdominios y tecnologías, así como la identificación preliminar de posibles vulnerabilidades.

La solución propuesta integra múltiples utilidades y técnicas ampliamente utilizadas en entornos profesionales de ciberseguridad (como Nmap, Feroxbuster, WhatWeb o WPScan), y se construye sobre una arquitectura modular y escalable que permite incorporar nuevas funcionalidades de forma sencilla en el futuro. A través de una interfaz web intuitiva, el usuario puede gestionar proyectos, ejecutar escaneos personalizados y consultar resultados de manera clara y estructurada.

Esta herramienta permite optimizar el flujo de trabajo durante las fases iniciales de una auditoría, automatizando tareas repetitivas, centralizando la información obtenida y reduciendo significativamente los tiempos de ejecución. Por otra parte, incorpora un sistema de inteligencia artificial basado en modelos generativos, que facilita la interacción con el contexto del escaneo para obtener recomendaciones específicas o definir próximos pasos según las debilidades detectadas. Además, genera automáticamente informes técnicos versionados, lo que simplifica la documentación del proceso y mejora su trazabilidad sin necesidad de intervención manual.

\noindent\textbf{Palabras clave}: \printkeywords.

\chapter*{Abstract}
This project presents the development of a tool aimed at automating fundamental tasks in security audits, with a specific focus on the reconnaissance phase. This stage, which is critical in any penetration testing process, includes activities such as passive and active information gathering, port and service scanning, subdomain and technology enumeration, and the preliminary identification of potential vulnerabilities.

The proposed solution integrates several tools and techniques widely used in professional cybersecurity environments (such as Nmap, Feroxbuster, WhatWeb, or WPScan), and is built on a modular and scalable architecture that facilitates the future incorporation of new functionalities. Through an intuitive web interface, users can efficiently manage projects, execute customized scans, and review results in a clear and structured way.

The tool helps optimize workflow during the initial stages of a security audit by automating repetitive tasks, centralizing collected data, and significantly reducing execution time. Moreover, it includes an artificial intelligence system based on generative models, which allows users to interact with the scan context to receive specific recommendations or define next steps based on detected weaknesses. Additionally, it automatically generates versioned technical reports, streamlining the documentation process and improving traceability without requiring manual intervention.

\noindent\textbf{Keywords}: Security, Ethical Hacking, Auditing, Automation, Tool.

% Índices con enlaces de las secciones, tablas y figuras
\newpage
\pagestyle{plain}
\pagenumbering{Roman}
\setcounter{page}{1}
\tableofcontents
%\pagebreak
\newpage
\listoftables
\newpage
\listoffigures
\newpage


%%%%%%% ESTRUCTURA DEL TFG %%%%%%%%%%%%%%
% A.PORTADA (según la estructura indicada a continuación) 
% B.CONTRAPORTADA (según la estructura indicada a continuación) 
% C.ÍNDICE GENERAL DE CONTENIDOS  
% D.ÍNDICE DE TABLAS  
% E.ÍNDICE DE FIGURAS 
% F.RESUMEN (Podrá incluirse también en inglés, si así lo indica el Tutor1) 
% G.CUERPO DEL TRABAJO (según la estr
% uctura indicada a continuación) 
% H.REFERENCIAS BIBLIOGRÁFICAS 
% (Según norma ISO690)
% I.ANEXOS, si los hubiera
%%%%%%%%%%%%%%%%%%%%%%%%%%%%%%%%%%%%%%%
%%%%%% Cuerpo del trabajo
% 1.INTRODUCCIÓN 
% 2.OBJETIVOS 
% 3.ANTECEDENTES / ESTADO DEL ARTE 
% 4.MÉTODOLOGÍA 
% 5.IMPLEMENTACIÓN Y DESARROLLO (Cuando proceda) 
% 6.RESULTADOS Y DISCUSIÓN 
% 7.CONCLUSIONES 
%%%%%%%%%%%%%%%%%%%%%%%%%%%%%%%%%%%%%%%%%%%
%%%%%%%%%%%%%%%%%%%%%%%%%%%%%%%%%%%%%%%%%%%
% Tamaño: Normalizado UNE A-4, salvo planos. 
% Tipo y tamaño de letra del texto: Times New Roman 12 pt, Arial 12 pt o similar. 
% Interlineado del texto: 1,5 líneas. 
% Márgenes del texto: Superior, Inferior y Derecha, 2,5 cm; Izquierda, 4 cm. 
% Numeración de páginas en margen inferior derecha y tamaño 8 pt. 
% Las  figuras  serán  numeradas  y  tituladas  debajo  de  las  mismas  (indicando  su  fuente  si  no  son  de  elaboración  
% propia). 
% Las  tablas  serán  numeradas  y  tituladas  encima  de  las  mismas  (indicando  su  fuente  si  no  son  de  elaboración  
% propia). 
%%%%%%%%%%%%%%%%%%%%%%%%%%%%%%%%%%%%%%%%%%%
%%%%%%%%%%%%%%%%%%%%%%%%%%%%%%%%%%%%%%%%%%%
%%%%%%%%%%%%%%%%%%%%%%%%%%%%%%%%%%%%%%%%%%%
%%%%%%%%%%%%%%%%%%%%%%%%%%%%%%%%%%%%%%%%%%%


\pagenumbering{arabic}
\pagestyle{fancy}
\setcounter{page}{1}

%%%%%%%% INTRODUCCIÓN %%%%%%%%
\import{chapters}{01_introduccion.tex}

%%%%%%%% OBJETIVOS %%%%%%%%
\import{chapters}{02_objetivos.tex}

%%%%%%%% ESTADO DEL ARTE %%%%%%%%
\import{chapters}{03_estado_del_arte.tex}

%%%%%%%% METODOLOGÍA %%%%%%%%
\import{chapters}{04_metodologia.tex}

%%%%%%%% IMPLEMENTACIÓN Y DESARROLLO %%%%%%%%
\import{chapters}{05_implementacion_y_desarrollo.tex}

%%%%%%%% RESULTADOS %%%%%%%%
\import{chapters}{06_resultados.tex}

%%%%%%%% CONCLUSIONES Y TRABAJOS FUTUROS %%%%%%%%
\import{chapters}{07_conclusiones_y_trabajos_futuros.tex}


% Si no se van a incorporar anexos, comentar el comando '\startappendices' junto a todos los imports de la carpeta 'appendices'
%%%%%%%%%%%%%%%%%%%%%%%%%%%%%%%%%%
%%%%%%%%%%% ANEXOS %%%%%%%%%%%%%%%
%%%%%%%%%%%%%%%%%%%%%%%%%%%%%%%%%%
%\startappendices

%%%%%%%% ANEXO 1 %%%%%%%%
%\import{appendices}{01_ejemplo_de_anexo.tex}

\pagebreak
%%%%%%%%%%%%%%%%%%%%%%%%%%%%%%%%%%
%%%%%%%%%% AL FINAL %%%%%%%%%%%%%%
%%%%%%%%%%%%%%%%%%%%%%%%%%%%%%%%%%
\pagestyle{empty}
%%%% https://en.wikibooks.org/wiki/LaTeX/Bibliography_Management
%%%% https://www.bibtex.com
%%%%%%%% BIBLIOGRAFÍA %%%%%%%%
% Se puede cambiar el estilo de la bibliografía, por defecto es 'unsrturl'
% Por ejemplo: \makebibliography{alphaurl}
\makebibliography
\end{document}